\clearpage
\newcounter{ch2checkpointappendix}
\renewcommand{\thech2checkpointappendix}{\thechapter.A}
\refstepcounter{ch2checkpointappendix}
\section*{Chapter 2 Appendix: Checkpoint-Selection Diagnostics}
\addcontentsline{toc}{section}{Chapter 2 Appendix: Checkpoint-Selection Diagnostics}
\label{app:checkpoint_selection}

This appendix records checkpoint-selection diagnostics for the three
paper-level post-trained models. \textit{Model chk-1} is the analytical SFT
checkpoint 200. \textit{Model chk-2} is the independent Minutes-rewriting SFT
branch represented by checkpoint 50 loaded as a LoRA adapter over the
exact-merged \textit{Model chk-1} checkpoint~200 parent. \textit{Model chk-3}
is the independent policy-decision branch, with SFT checkpoint 38 as its
pre-GRPO state and GRPO checkpoint 450 as a post-GRPO exploratory evaluation
state. These paper names are task labels and should not be inferred from older
repository-side \texttt{chkN} identifiers.

The stages use different prompts, samples, and admission rules. The
\textit{Model chk-1} gate uses a frozen eight-case generation-stability probe;
the post-checkpoint-200 neighbour sweep reported below is retrospective. The
\textit{Model chk-2} gate uses eight deterministically selected validation rows
from the 42-row monitoring split and combines zero-tolerance generation checks
with tolerant checkpoint-level thresholds over source fidelity and two fresh
validators. It is selection-only, uses official pre-action Minutes in Validator
B, and is not a leakage-safe population evaluation. The \textit{Model chk-3}
SFT qualification probe is train-stratified, whereas its GRPO gate uses separate
selection, blind-validation, and train-only retention panels. Case counts and
pass rates are therefore not comparable across model stages.

Table~\ref{tab:app:chk1_post_cp200_hard_gates} audits late-trajectory
free-running analysis generation, and
Table~\ref{tab:app:chk2_cp40_cp60_selection} reports the current
\textit{Model chk-2} candidate comparison. Tables~\ref{tab:app:chk3_sft_qualification_gates}
and~\ref{tab:app:chk3_grpo_selection_gates} distinguish the decision-SFT
qualification gate from the subsequent GRPO checkpoint-selection gate, while
Table~\ref{tab:app:chk3_sft_cp38_cp39_selection} records why SFT checkpoint 38,
rather than the terminal checkpoint 39, initializes GRPO.

\subsection*{Learning-Rate Stability Diagnostics for Model chk-1}
\label{app:chk1_lr_selection}

The peak learning rate for \textit{Model chk-1} was chosen through staged
generation-stability diagnostics rather than an exhaustive hyperparameter
search. The first controlled smoke comparison held the dataset, sample order,
random seeds, LoRA configuration, optimizer, warm-up schedule, batch size, and
generation prompts fixed while reducing the peak learning rate from
\(5\times10^{-5}\) to \(1\times10^{-5}\). At \(5\times10^{-5}\), the first
matched held-out completion reached the 3,072-token limit, failed to emit both
an end-of-sequence token and a valid reasoning boundary, and recorded full-text
and tail four-gram repetition rates of 0.9231 and 0.9647, respectively. This
single case satisfied the preregistered fail-fast condition, so the remaining
seven cases were not generated. By contrast, the ten-step
\(1\times10^{-5}\) smoke checkpoint passed all eight fixed probe cases.

Passing the short smoke test did not establish stability over a complete
training trajectory. The three-epoch \(1\times10^{-5}\) run completed all 255
optimizer updates and attained a terminal validation loss of 1.300019, but its
terminal checkpoint failed the same generation gate: one of eight completions
reached the token limit, only seven satisfied the output contract, and three
were classified as catastrophically repetitive. The lower validation loss
therefore did not imply reliable free-running generation. The peak learning
rate was subsequently reduced to \(1\times10^{-6}\), and a new three-epoch run
was initialized from \textit{Model chk-0}. In that run, checkpoints 80, 150,
200, and 255 all passed the fixed eight-case generation probe. Checkpoint 200
was retained as the paper-facing \textit{Model chk-1} because it combined the
minimum logged validation loss (1.686099) with a fully passing generation
probe; checkpoint 255 had an essentially unchanged terminal validation loss
of 1.686131 and also passed, but exhibited slightly greater mean four-gram
repetition (0.140996 versus 0.134664).

\begin{table}[htbp]
    \centering
    \scriptsize
    \setlength{\tabcolsep}{3.5pt}
    \renewcommand{\arraystretch}{1.15}
    \begin{tabularx}{\textwidth}{@{}l>{\raggedright\arraybackslash}Xcccccc@{}}
        \toprule
        \makecell{\textbf{Peak learning}\\\textbf{rate}} &
        \textbf{Diagnostic state} &
        \textbf{Cases} &
        \makecell{\textbf{Contract}\\\textbf{valid}} &
        \textbf{Capped} &
        \makecell{\textbf{Catastrophic}\\\textbf{repetition}} &
        \makecell{\textbf{Mean full-text}\\\textbf{4-gram repetition}} &
        \textbf{Gate} \\
        \midrule
        \(5\times10^{-5}\) & ten-step smoke, cp10 & 1\textsuperscript{a} & 0/1 & 1/1 & 1/1 & 0.923102 & Failed \\
        \(1\times10^{-5}\) & ten-step smoke, cp10 & 8 & 8/8 & 0/8 & 0/8 & 0.094759 & Passed \\
        \(1\times10^{-5}\) & three-epoch run, cp255 & 8 & 7/8 & 1/8 & 3/8 & 0.408608 & Failed \\
        \(1\times10^{-6}\) & three-epoch run, cp200 & 8 & 8/8 & 0/8 & 0/8 & 0.134664 & Passed \\
        \bottomrule
    \end{tabularx}
    \begin{minipage}{0.96\textwidth}
        \vspace{2pt}
        \footnotesize
        \textit{Note:} \textsuperscript{a}\ The \(5\times10^{-5}\) probe stopped
        after its first case triggered the fail-fast rule; its row is therefore
        a failure demonstration, not an estimate of an eight-case failure rate.
        ``Catastrophic repetition'' follows the frozen probe thresholds of at
        least 0.50 full-text or 0.60 tail token four-gram repetition. All rows
        use the same frozen probe manifest, but they represent an adaptive,
        staged diagnostic sequence rather than a balanced or statistically
        powered learning-rate sweep.
    \end{minipage}
    \caption{Staged learning-rate stability diagnostics for \textit{Model
    chk-1}. Selection prioritized valid and non-degenerate free-running
    generation; validation loss alone was not treated as sufficient evidence
    of model reliability.}
    \label{tab:app:chk1_lr_stability}
\end{table}

\subsection*{Post-cp200 Generation-Stability Replay for Model chk-1}

To determine whether the retained checkpoint was uniquely stable or whether
stability persisted across the late loss plateau, checkpoints 210--250 were
replayed retrospectively on 23 August 2026 using the same frozen eight-case
manifest, seeds, decoding settings, \textit{Model chk-0} baseline, and runtime
versions as the original checkpoint-200 and checkpoint-255 probes. These
intermediate replays did not participate in the original checkpoint-selection
decision. Table~\ref{tab:app:chk1_post_cp200_hard_gates} therefore provides a
matched descriptive audit of the late trajectory rather than new prospective
selection evidence.

\begin{table}[htbp]
    \centering
    \scriptsize
    \setlength{\tabcolsep}{3.25pt}
    \renewcommand{\arraystretch}{1.15}
    \begin{tabular}{@{}lcccccc@{}}
        \toprule
        \textbf{Checkpoint} &
        \makecell{\textbf{Training / evaluation}\\\textbf{loss}} &
        \makecell{\textbf{Finite / EOS /}\\\textbf{contract-valid}} &
        \makecell{\textbf{Cap / catastrophic /}\\\textbf{periodic}} &
        \makecell{\textbf{Mean full-text}\\\textbf{4-gram repetition}} &
        \makecell{\textbf{$\Delta$ versus chk-0}\\\textbf{mean / maximum}} &
        \textbf{Gate} \\
        \midrule
        \texttt{cp200} (retained) & 1.729560 / 1.686099 & 8 / 8 / 8 & 0 / 0 / 0 & 0.134664 & 0.051945 / 0.144158 & Passed \\
        \texttt{cp210}\textsuperscript{a} & 1.694860 / 1.686223 & 8 / 8 / 8 & 0 / 0 / 0 & 0.148501 & 0.065781 / 0.135519 & Passed \\
        \texttt{cp220}\textsuperscript{a} & 1.709530 / 1.686188 & 8 / 8 / 8 & 0 / 0 / 0 & 0.142607 & 0.059888 / 0.161819 & Passed \\
        \texttt{cp230}\textsuperscript{a} & 1.724720 / 1.686177 & 8 / 8 / 8 & 0 / 0 / 0 & 0.116562 & 0.033842 / 0.158795 & Passed \\
        \texttt{cp240}\textsuperscript{a} & 1.704090 / 1.686112 & 8 / 8 / 8 & 0 / 0 / 0 & 0.111440 & 0.028720 / 0.146266 & Passed \\
        \texttt{cp250}\textsuperscript{a} & 1.705740 / 1.686216 & 8 / 8 / 8 & 0 / 0 / 0 & 0.150136 & 0.067417 / 0.141370 & Passed \\
        \texttt{cp255} (terminal) & 1.722860 / 1.686131 & 8 / 8 / 8 & 0 / 0 / 0 & 0.140996 & 0.058276 / 0.171121 & Passed \\
        \bottomrule
    \end{tabular}
    \begin{minipage}{0.97\textwidth}
        \vspace{2pt}
        \footnotesize
        \textit{Note:} \textsuperscript{a}\ Retrospective matched replay; the
        original selection record contained probes for checkpoints 200 and 255
        but not for these intermediate checkpoints. Each three-part count is
        the number of cases, out of eight, in the order stated in its column
        heading. Each probe contains six unique prompts plus two additional
        greedy decodes. The
        frozen comparator requires finite
        output for every case, no completion at the 3,072-token cap, no strict
        periodic tail, no catastrophic repetition (full-text four-gram rate
        $\geq 0.50$ or tail rate $\geq 0.60$), no contract-validity loss relative
        to the fully valid chk-0 baseline, a mean repetition increase no greater
        than 0.10, and a per-case increase no greater than 0.20. EOS is reported
        as a diagnostic but is not a separate comparator threshold. In the loss
        column, the training value is the trailing mean of the ten step-level
        losses ending at the reported checkpoint,
        $\overline{L}^{\mathrm{train}}_{c,10}=10^{-1}\sum_{j=0}^{9}
        L^{\mathrm{train}}_{c-j}$, followed by evaluation loss. Neither loss is
        part of the generation gate; the checkpoint-255 evaluation value is the
        terminal evaluation loss. Full-text and 1,024-token-tail repetition
        values coincide in every row because no retained completion contains
        more than 1,023 content tokens.
    \end{minipage}
    \caption{Post-cp200 generation-stability hard-gate comparison for
    \textit{Model chk-1}. All evaluated checkpoints pass the frozen gate, so
    the table does not establish checkpoint 200 as uniquely stable. Its
    retention remains supported by the minimum logged evaluation loss together
    with a passing contemporaneous generation probe.}
    \label{tab:app:chk1_post_cp200_hard_gates}
\end{table}

\subsection*{Tolerant Checkpoint-Level Selection for Model chk-2}

The current Minutes-rewriting branch starts from the exact-merged
\textit{Model chk-1} checkpoint~200 parent and trains a separate LoRA adapter on
305 training rows for three epochs, producing 60 optimizer steps. The terminal
root adapter is not selected automatically. Instead, checkpoints 40, 50, and
60 are loaded as adapters over the same frozen parent and compared on a fixed
screening panel. Later paper evaluations use the selected checkpoint~50 adapter
in this same parent-plus-adapter form; this comparison is not an exact-merged
candidate replay.

The panel contains eight rows drawn deterministically from the 42-row
validation split without access to the supervised target. Five rows represent
fixed prompt-length quantiles; the remaining three prioritize high numerical
density, high date density, and attribution content, with distinct meetings
preferred throughout. The realized panel covers eight meetings, seven atomic
topics, and both section-style strata. Every candidate uses the same prompts
under greedy decoding, a 2,048-token completion limit, four-bit NF4 inference,
and the same recorded seed. No test row is generated during selection.

The selection contract deliberately separates strict per-row verdicts from
tolerant checkpoint-level admission. For row \(i\) and candidate \(c\), let
\(Z_{i,c}\) indicate passage of the zero-tolerance structure and degeneration
checks. These checks require a finite nonempty completion, terminal EOS,
compliance with the generation and 4,096-token sequence limits, exactly one
reasoning boundary, nonempty native reasoning, one 20--400-word final
paragraph, no prohibited headings, lists, JSON, citations, credentials, or
control markers, no strict periodic tail, and full-text and tail token
four-gram repetition below 0.50 and 0.60, respectively. Let \(D_{i,c}\)
indicate deterministic source fidelity: the normalized numeric multiset and
the complete date and attribution sets must be preserved, and the final
paragraph must not be an exact copy of the full source analysis. Finally, let
\(G_{i,c}\) indicate passage of the complete local deterministic gate, which
also rejects a duplicated final sentence before either validator is called.

Validator A receives the source analysis, student reasoning, and rewritten
paragraph, but not the official Minutes, and checks bidirectional claim support
and reasoning--answer compatibility. Validator B receives the rewritten
paragraph and the corresponding meeting's pre-action official Minutes, but not
the source analysis or reasoning, and applies the six-dimensional style
contract. Denote their strict binary verdicts by \(A_{i,c}\) and \(B_{i,c}\),
and let \(s^B_{i,c}\) be Validator B's mean style score. A row passes the full
chain only when

\[
    H_{i,c}=G_{i,c}A_{i,c}B_{i,c}=1.
\]

A candidate is admitted only if all of the following checkpoint-level
requirements hold simultaneously:

\[
\begin{gathered}
    \sum_{i=1}^{8}(1-Z_{i,c})=0,
    \qquad \overline{D}_{c}\geq 0.75,
    \qquad \overline{A}_{G,c}\geq 0.75, \\
    \overline{B}_{A,c}\geq 0.50,
    \qquad \overline{s}^{B}_{A,c}\geq 6.0,
    \qquad \overline{H}_{c}\geq 0.25,
    \qquad U_c=0,
\end{gathered}
\]

where the Validator-A rate is conditional on complete deterministic-gate passage,
the Validator-B rate and mean score are conditional on Validator-A passage,
and \(U_c\) is the number of unresolved required validator calls. The mean
Validator-B score includes every Validator-A-pass case with a valid score,
including cases that do not receive a Validator-B PASS verdict. If more than
one candidate is eligible, the receipt-defined ranking order is full-chain pass rate,
Validator-B mean score, Validator-A pass rate, deterministic-fidelity pass
rate, and then the earlier checkpoint step.

\begin{table}[htbp]
    \centering
    \scriptsize
    \setlength{\tabcolsep}{2.3pt}
    \renewcommand{\arraystretch}{1.15}
    \begin{tabularx}{\textwidth}{@{}>{\raggedright\arraybackslash}Xcccccccc@{}}
        \toprule
        \textbf{Candidate} &
        \makecell{\textbf{Evaluation}\\\textbf{loss}} &
        \makecell{\textbf{\(Z\)}\\\textbf{failures}} &
        \makecell{\textbf{\(D\) pass}\\\textbf{(of 8)}} &
        \makecell{\textbf{\(A\) pass}\\\textbf{(of \(G\))}} &
        \makecell{\textbf{\(B\) pass}\\\textbf{(of \(A\))}} &
        \makecell{\textbf{Mean}\\\textbf{\(B\) score}} &
        \makecell{\textbf{Full chain}\\\textbf{\(H\) (of 8)}} &
        \textbf{Admission} \\
        \midrule
        \texttt{cp40} & 1.529542 & 0 & 7/8 & 5/7 & 0/5 & 4.3333 & 0/8 & Failed \\
        \texttt{cp50} (retained) & 1.526579 & 0 & 6/8 & 5/6 & 3/5 & 6.7333 & 3/8 & \textbf{Passed} \\
        \texttt{cp60} (terminal) & \textbf{1.526236} & 0 & 6/8 & 5/6 & 1/5 & 6.3333 & 1/8 & Failed \\
        \bottomrule
    \end{tabularx}
    \begin{minipage}{0.97\textwidth}
        \vspace{2pt}
        \footnotesize
        \textit{Note:} Evaluation loss is reported as an optimization
        diagnostic and is not an admission criterion. All three candidates have
        \texttt{strict\_all\_rows\_pass=false}. For every candidate shown,
        the \(G\)-pass count equals the reported \(D\)-pass count because there
        are no \(Z\) failures or duplicated-final-sentence failures. The
        checkpoint~60 row uses the
        final remediated record: one initially unresolved Validator-B report was
        resolved by contract-only remediation without changing the generated
        candidate, leaving zero unresolved calls. ``Retained'' identifies the
        paper-facing adapter checkpoint. The selection receipt is selection-only
        and does not itself authorize a merge or canonical-DAG promotion.
    \end{minipage}
    \caption{Checkpoint-level candidate comparison for the current
    \textit{Model chk-2} Minutes-rewriting branch. Checkpoint 50 is the only
    candidate satisfying every tolerant aggregate admission requirement; this
    does not mean that every checkpoint~50 row passes the strict full chain.}
    \label{tab:app:chk2_cp40_cp60_selection}
\end{table}

Checkpoint 40 passes the zero-tolerance and aggregate deterministic-fidelity
requirements but falls below the Validator-A, Validator-B, mean-style-score,
and full-chain thresholds. Checkpoint 60 passes the zero-tolerance,
deterministic-fidelity, Validator-A, and mean-style-score requirements, but its
Validator-B pass rate is \(1/5=0.20\) and its full-chain pass rate is
\(1/8=0.125\). Checkpoint 50 is the sole admitted candidate: six of eight rows
pass deterministic fidelity, five of those six pass Validator A, three of
those five pass Validator B, and three of eight pass the complete chain.

The terminal checkpoint has the minimum monitoring loss, 1.5262358, compared
with 1.5265795 at checkpoint 50. The difference is only 0.0003437, or 0.0225\%
relative to the checkpoint~50 value, and loss is not part of the quality gate.
Checkpoint 50 is therefore retained because it is the only candidate to pass
the documented checkpoint-level contract, not because it minimizes validation
loss and not because it passes every sample-level hard gate.

This eight-row panel provides internal selection evidence only. It is neither
an untouched test nor a basis for population inference: its rows come from the
monitoring split, and Validator B directly uses corresponding official
pre-action Minutes as a style comparator. BERTScore, MPNet similarity, Core8
interventions, and historical association results are downstream diagnostics
of the already selected checkpoint and do not participate in this selection
rule. The supervised reference paragraphs remain synthetic teacher rewrites;
official Minutes are used by Validator B only for the stated style comparison.

\clearpage
\subsection*{Separate SFT and GRPO Hard Gates for Model chk-3}

Paper-level \textit{Model chk-3} corresponds to repository artifact
\texttt{chk4}. Its two training stages answer different gate questions. The
decision-SFT gate asks whether an SFT candidate is sufficiently well behaved to
initialize GRPO; the matched replay below distinguishes checkpoint 38 from the
terminal checkpoint 39. The GRPO gate asks whether any later policy checkpoint
meets the frozen checkpoint-selection requirements. Passing the first gate does
not imply passing the second.

\begin{table}[htbp]
    \centering
    \scriptsize
    \setlength{\tabcolsep}{3.5pt}
    \renewcommand{\arraystretch}{1.15}
    \begin{tabularx}{\textwidth}{@{}>{\raggedright\arraybackslash}Xccc>{\raggedright\arraybackslash}X@{}}
        \toprule
        \textbf{Gate component} &
        \textbf{Frozen requirement} &
        \makecell{\textbf{Observed at}\\\textbf{SFT cp38}} &
        \textbf{Result} &
        \textbf{Function of the gate} \\
        \midrule
        Reasoning-boundary delivery & Rate $\geq 0.75$ & 14/15 = 0.9333 & Passed & Requires a usable native reasoning boundary before the terminal action. \\
        Completion cap & Rate $\leq 0.25$ & 1/15 = 0.0667 & Passed & Limits truncated generations under the 1,024-token probe ceiling. \\
        Strict periodic tail & Count $=0$ & 0/15 & Passed & Rejects mechanically looping completions. \\
        Sampled direction correctness by class & At least one per class & \makecell{hold 3; hike 1; cut 1} & Passed & Requires non-zero action signal for each policy direction. \\
        Sampled non-zero reward by class & At least one per class & \makecell{hold 3; hike 1; cut 2} & Passed & Prevents an all-zero reward group. \\
        Sampled reward dispersion by class & Standard deviation $>0$ & \makecell{hold 0.4330; hike 0.1028;\\cut 0.3074} & Passed & Ensures within-group reward variation for GRPO. \\
        \bottomrule
    \end{tabularx}
    \begin{minipage}{0.97\textwidth}
        \vspace{2pt}
        \footnotesize
        \textit{Note:} The probe contains three direction-stratified groups,
        each with one greedy and four sampled generations. Class-specific gates
        use the four sampled generations. Exact magnitude is reported elsewhere
        as a diagnostic but is not an SFT qualification gate, because the
        validation-only 50/75-basis-point hike magnitudes are absent from the
        sealed training split. Checkpoint 38 passes this qualification gate and
        is exact-merged as the pre-GRPO SFT state of \textit{Model chk-3}.
    \end{minipage}
    \caption{Decision-SFT qualification hard gates for \textit{Model chk-3}
    at checkpoint 38. These gates determine whether the SFT state may initialize
    GRPO; they do not establish held-out decision quality.}
    \label{tab:app:chk3_sft_qualification_gates}
\end{table}

\subsubsection*{Why Checkpoint 38 Rather than Checkpoint 39?}

The terminal SFT checkpoint is not automatically preferred to an earlier
checkpoint. Checkpoints 38 and 39 were evaluated with the same
direction-stratified qualification probe and the same frozen thresholds.
Table~\ref{tab:app:chk3_sft_cp38_cp39_selection} separates the trainer-reported
loss observations from the behavioral quantities that actually determine
qualification.

\begin{table}[htbp]
    \centering
    \scriptsize
    \setlength{\tabcolsep}{3.5pt}
    \renewcommand{\arraystretch}{1.15}
    \begin{tabularx}{\textwidth}{
        @{}
        >{\raggedright\arraybackslash}X
        >{\raggedright\arraybackslash}X
        >{\centering\arraybackslash}p{0.22\textwidth}
        >{\centering\arraybackslash}p{0.22\textwidth}
        @{}
    }
        \toprule
        \textbf{Diagnostic} &
        \textbf{Frozen requirement} &
        \textbf{SFT cp38} &
        \textbf{SFT cp39} \\
        \midrule
        Per-step training loss
        & Not a qualification gate
        & 1.6149
        & 1.3127 \\
        Scheduled evaluation loss
        & Not a qualification gate
        & Not recorded at step 38
        & 1.601505 \\
        Reasoning-boundary delivery
        & Overall rate $\geq 0.75$
        & $14/15=0.9333$
        & $13/15=0.8667$ \\
        Completion cap
        & Overall rate $\leq 0.25$
        & $1/15=0.0667$
        & $2/15=0.1333$ \\
        Strict periodic tail
        & Count $=0$
        & 0
        & 0 \\
        \makecell[l]{Sampled direction-correct\\coverage: hold / hike / cut}
        & At least one per class
        & $3 / 1 / 1$
        & $2 / 1 / \mathbf{0}$ \\
        \makecell[l]{Sampled non-zero-reward\\coverage: hold / hike / cut}
        & At least one per class
        & $3 / 1 / 2$
        & $3 / 2 / 3$ \\
        \makecell[l]{Sampled reward standard deviation:\\hold / hike / cut}
        & Greater than zero for every class
        & $0.4330 / 0.1028 / 0.3074$
        & $0.4878 / 0.4263 / 0.0217$ \\
        Overall qualification
        & Every behavioral gate must pass
        & \textbf{Passed}
        & \makecell{\textbf{Failed}: no sampled\\cut-direction success} \\
        \bottomrule
    \end{tabularx}
    \begin{minipage}{0.97\textwidth}
        \vspace{2pt}
        \footnotesize
        \textit{Note:} Each checkpoint is evaluated on one hold, one hike, and
        one cut training prompt, with one greedy and four sampled generations
        per prompt. The class-specific gates use the four sampled generations.
        Checkpoint 39 produces a correct cut under greedy decoding but produces
        no correct cut direction in its four sampled generations. Training loss
        is a batch-level statistic, and no scheduled evaluation was performed at
        step 38; neither loss quantity is part of this behavioral gate.
    \end{minipage}
    \caption{Matched SFT qualification comparison between checkpoints 38 and
    39 for \textit{Model chk-3}. Checkpoint 39 is the terminal SFT state but
    fails the frozen sampled cut-direction coverage requirement.}
    \label{tab:app:chk3_sft_cp38_cp39_selection}
\end{table}

Although checkpoint 39 records the minimum logged SFT evaluation loss, that
value is not a matched checkpoint-38 comparison because evaluation was
scheduled only at steps 13, 26, and 39. More importantly, loss is not the
selection criterion for GRPO initialization. Checkpoint 39 fails the frozen
requirement of at least one sampled direction-correct output in every policy
class: none of its four sampled cut completions predicts the correct direction.
Checkpoint 38 satisfies every qualification component and is therefore retained
and exact-merged as the pre-GRPO state of \textit{Model chk-3}. This small,
train-stratified probe establishes readiness to initialize GRPO, not population
superiority or a statistically unique loss optimum.

\begin{table}[htbp]
    \centering
    \scriptsize
    \setlength{\tabcolsep}{3.5pt}
    \renewcommand{\arraystretch}{1.15}
    \begin{tabular}{@{}lccccccc@{}}
        \toprule
        \textbf{GRPO checkpoint} &
        \makecell{\textbf{Mean}\\\textbf{reward}} &
        \makecell{\textbf{Direction}\\\textbf{accuracy}} &
        \makecell{\textbf{Exact-action}\\\textbf{accuracy}} &
        \makecell{\textbf{Delivery-valid}\\\textbf{rate}} &
        \textbf{Cap rate} &
        \makecell{\textbf{Periodic-tail}\\\textbf{count}} &
        \textbf{Gate} \\
        \midrule
        \texttt{cp310} & 0.227083 & 0.277778 & 0.222222 & 0.444444 & 0.111111 & 0 & Failed \\
        \texttt{cp360} & 0.106944 & 0.222222 & 0.166667 & 0.388889 & 0.000000 & 0 & Failed \\
        \texttt{cp410} & 0.093750 & 0.111111 & 0.111111 & 0.555556 & 0.055556 & 0 & Failed \\
        \texttt{cp440}\textsuperscript{a} & 0.122917 & 0.166667 & 0.111111 & 0.388889 & 0.166667 & 0 & Failed \\
        \texttt{cp450}\textsuperscript{b} & 0.143750 & 0.222222 & 0.166667 & 0.333333 & 0.166667 & 0 & Failed \\
        \texttt{cp460}\textsuperscript{a} & 0.078819 & 0.111111 & 0.055556 & 0.333333 & 0.111111 & 0 & Failed \\
        \texttt{cp468} (terminal) & 0.069444 & 0.055556 & 0.055556 & 0.333333 & 0.166667 & 0 & Failed \\
        \bottomrule
    \end{tabular}
    \begin{minipage}{0.97\textwidth}
        \vspace{2pt}
        \footnotesize
        \textit{Note:} The selection-panel eligibility gate requires a
        delivery-valid rate of at least 0.75, a cap rate no greater than 0.25,
        and zero periodic-tail cases. The reported panel contains 18 completions
        from nine selection-panel validation prompts; it contains no cut example.
        Every checkpoint fails the delivery threshold. \textsuperscript{a}
        Checkpoints 440 and 460 are post-hoc local-neighbour probes.
        \textsuperscript{b}\ Checkpoint 450 is the best-observed exploratory
        candidate, not a selected checkpoint. On its blind follow-up it records
        reward 0.1725, direction accuracy 0.30, exact-action accuracy 0.15,
        delivery validity 0.30, cap rate 0.10, and zero periodic tails. It covers
        cut, hold, and hike on the train-only retention panel, but fails the
        blind delivery gate. The authoritative receipt therefore records
        \texttt{best\_observed\_checkpoint\_step=450},
        \texttt{selected\_checkpoint\_step=null}, and
        \texttt{provisional\_no\_candidate\_passed\_all\_gates}. The later
        13-meeting test likewise records \texttt{quality\_not\_demonstrated}.
    \end{minipage}
    \caption{GRPO checkpoint-selection hard gates for \textit{Model chk-3}.
    No checkpoint satisfies the frozen eligibility rule; checkpoint 450 was
    evaluated only as the best-observed exploratory candidate through an
    explicit override.}
    \label{tab:app:chk3_grpo_selection_gates}
\end{table}
